%\documentclass[12pt,oneside]{amsart}
\documentclass{scrartcl}

%\documentclass{article}

\usepackage{amsthm,amsmath,amssymb}
\usepackage[numbers]{natbib}

%\usepackage[colorlinks]{hyperref}
\usepackage{hyperref}
\hypersetup{
    colorlinks,%
    citecolor=black,%
    filecolor=black,%
    linkcolor=black,%
    urlcolor=black
}
\usepackage{hypernat}
\usepackage[top=2.5cm, bottom=2.5cm, left=2.8cm,
right=2.8cm]{geometry}
\usepackage{graphicx}


%\setlength{\parskip}{0pt}
%\setlength{\parsep}{0pt}
%\setlength{\headsep}{0pt}
%\setlength{\topskip}{0pt}
%\setlength{\topmargin}{0pt}
%\setlength{\topsep}{0pt}
%\setlength{\partopsep}{0pt}

\linespread{1}

%\usepackage{marvosym}

%\textwidth = 6.5 in \textheight = 9.2 in \oddsidemargin = 0.0 in
%\evensidemargin = 0.0 in \topmargin = -0.0 in \headheight = 0.0 in
%\headsep = 0.0in \parskip = 0.0in \parindent = 0.0in
%\def\baselinestretch{0.871}


\newtheorem{theorem}{Theorem}[section]
\newtheorem{proposition}[theorem]{Proposition}
\newtheorem{problem}[theorem]{Problem}
\newtheorem{fact}[theorem]{Fact}
\newtheorem{lemma}[theorem]{Lemma}
\newtheorem{defn}[theorem]{Definition}
\newtheorem{corollary}[theorem]{Corollary}
\newtheorem{remark}{Remark}[section]
\newtheorem{example}[theorem]{Example}

\newcommand{\E}{{\mathbb E}}
\newcommand{\se}{{\mathcal{E}}}
\newcommand{\R}{{\mathbb R}}
\renewcommand{\P}{{\mathbb P}}
\newcommand{\ac}{{\mathcal{A}}}
\newcommand{\A}{{\mathcal{A}}}
\newcommand{\B}{{\mathcal{B}}}
\newcommand{\C}{{\mathcal{C}}}
\newcommand{\M}{{\mathcal{M}}}
\newcommand{\Mf}{{\mathfrak{M}}}
\newcommand{\T}{{\mathcal{T}}}
\newcommand{\Z}{{\mathcal{Z}}}
\newcommand{\K}{{\mathcal{K}}}
\newcommand{\lc}{{\mathcal{L}}}
\newcommand{\Ps}{{\mathcal{P}}}
\newcommand{\G}{{\mathcal{G}}}
\newcommand{\pa}{{\mathring{p}}}
\newcommand{\F}{{\cal F}}
\newcommand{\samp}{{\mathcal{X}}}
\newcommand{\X}{{\mathcal{X}}}
\newcommand{\hi}{{\hat{\Phi}}}
\newcommand{\data}{{\text{data}}}

\newcounter{rcnt}[section]
\renewcommand{\thercnt}{(\roman{rcnt})}

\def\qt#1{\qquad\text{#1}}

\def\argmin{\mathop{\rm argmin}}
\def\argmax{\mathop{\rm argmax}}


\setlength{\parskip}{1.4 \medskipamount}
%\sloppy
%\linespread{1.3}

\begin{document}

\title{STAT 238 - Bayesian Statistics  \\
  Lecture Fifteen} 
\subtitle{Spring 2026, UC Berkeley}

\author{Aditya Guntuboyina}


\date{25 Feb 2026}

\maketitle

\section{Linear Regression}

Our next topic is (multiple) linear regression. Here, we have one
response variable $y$ and 
$m$ covariates $x_1, \dots, x_m$  ($m = 1$ corresponds to simple linear
regression). We observe data on $n$ instances or subjects for all
these variables: $(y_i, x_{i1}, \dots, x_{im})$ for $i = 1, \dots,
n$. The multiple linear regression model (with normal errors) is given
by: 
\begin{equation}\label{lm}
  y_i = \beta_0 + \beta_1 x_{i1} + \dots + \beta_m x_{im} + \epsilon_i
~~  \text{with $\epsilon_i \overset{\text{i.i.d}}{\sim} N(0, \sigma^2)$}. 
\end{equation}
The default prior for linear regression is: 
\begin{equation*}
  \beta_0, \beta_1, \dots, \beta_m, \log \sigma
  \overset{\text{i.i.d}}{\sim} \text{unif}(-C, C)
\end{equation*}
for a very large positive $C$. The joint posterior density of
$\beta_0, \dots, \beta_m, \sigma$ is then given by
\begin{align*}
&  f_{\beta_0, \beta_1, \dots, \beta_m, \sigma \mid \text{data}}(\beta_0, \beta_1,
  \dots, \beta_m,
  \sigma) \\ &\propto  \sigma^{-n-1} \exp 
                                        \left(-\frac{S(\beta_0, \dots,
               \beta_m)}{2 \sigma^2}
    \right) I \left\{-C < \beta_0,
    \beta_1, \dots, \beta_m, \log \sigma < C \right\}. 
\end{align*}
where we use the notation
\begin{equation*}
  S(\beta_0, \beta_1, \dots, \beta_m)  := \sum_{i=1}^n (y_i - \beta_0 - \beta_1
      x_{i1} - \dots - \beta_m x_{im})^2
\end{equation*}
for the sum of squares. 

The 
posterior over only the coefficient parameters $\beta_0, \dots, \beta_m$ can
be obtained by integrating (or marginalizing) the parameter $\sigma$. 
\begin{align*}
 & f_{\beta_0, \beta_1, \dots, \beta_m \mid \text{data}} (\beta_0, \beta_1, \dots,
   \beta_m) \\  &= \int 
                                                              f_{\beta_0, 
                                                              \beta_1,
                  \dots, \beta_m,
                                                              \sigma
                                                              \mid
                                                              \text{data}}(\beta_0,
                                                              \beta_1,
                  \dots, \beta_m,
                                                              \sigma)
                                                              d\sigma
  \\
  &\propto I\{-C < \beta_0, \beta_1, \dots, \beta_m < C\}
    \int_{e^{-C}}^{e^C} \sigma^{-n-1} \exp 
                                        \left(-\frac{S(\beta_0, \dots,
    \beta_m)}{2 \sigma^2} \right) d\sigma \\
&\approx 
    I\{-C < \beta_0, \beta_1, \dots, \beta_m < C\} \int_{0}^{\infty} \sigma^{-n-1}
    \exp 
                                        \left(-\frac{S(\beta_0,\dots, \beta_m)}{2 \sigma^2}                               
  \right)  d\sigma \\
&=    I\{-C < \beta_0, \beta_1, \dots, \beta_m < C\}
  \left(\frac{1}{S(\beta_0, \dots, \beta_m)}\right)^{n/2} \int_{0}^{\infty} s^{-n-1} \exp
  \left(-\frac{1}{2s^2} \right) ds \\
&\propto   I\{-C < \beta_0, \beta_1, \dots, \beta_m < C\}
  \left(\frac{1}{S(\beta_0, \dots, \beta_m)}\right)^{n/2} \\
 &\approx   \left(\frac{1}{S(\beta_0, \dots, \beta_m)}\right)^{n/2} \propto 
  \left(\frac{S(\hat{\beta}_0, \hat{\beta}_1, \dots,
      \hat{\beta}_m)}{S(\beta_0, \beta_1, \dots, \beta_m)} 
  \right)^{n/2}
\end{align*}
where $\hat{\beta}_0, \dots, \hat{\beta}_m$ denote the least squares
estimators of $\beta_0, \dots, \beta_m$ (i.e., $(\hat{\beta}_0, \dots,
\hat{\beta}_m)$ minimizes $S(\beta_0, \dots, \beta_m)$ over all values
of $\beta_0, \dots, \beta_m$). 

Our posterior density for $\beta_0, \dots, \beta_m$ is thus:  
\begin{equation}\label{bayespost_lm}
   f_{\beta_0, \beta_1, \dots, \beta_m \mid \text{data}} (\beta_0, \beta_1, \dots,
   \beta_m) \propto \left(\frac{S(\hat{\beta}_0, \hat{\beta}_1, \dots,
      \hat{\beta}_m)}{S(\beta_0, \beta_1, \dots, \beta_m)} 
  \right)^{n/2}. 
\end{equation}
It turns out that this is a multivariate $t$-density. Here is some
basic background on multivariate $t$-densities.

\section{$t$-density}
The formula for the density corresponding to the $t$-distribution
$t_{p}(\mu, \Sigma, \nu)$ is (see
(\url{https://en.wikipedia.org/wiki/Multivariate_t-distribution}): 
\begin{align}
 f(x) &:= \frac{\Gamma((\nu + p)/2)}{\Gamma(\nu/2) \nu^{p/2} \pi^{p/2}
   \sqrt{\text{det}~\Sigma}} \left[\frac{1}{1 + \frac{1}{\nu} (x - \mu)^T
      \Sigma^{-1} (x - \mu)} \right]^{(\nu + p)/2} \nonumber \\
      &\propto  \left[\frac{1}{1 + \frac{1}{\nu} (x - \mu)^T 
      \Sigma^{-1} (x - \mu)} \right]^{(\nu + p)/2}. \label{tdens}
\end{align}
Here:
\begin{enumerate}
\item $p$ denotes dimension of the vector $x$ (this is a $p$-variate
  joint density)
\item $\mu$ is a $p \times 1$ vector called the location
\item $\Sigma$ is a $p \times p$ matrix called the scale matrix
\item $\nu > 0$   denotes the degrees of freedom. 
\end{enumerate}
Here is some more information about the $t$-density \eqref{tdens}:

\begin{enumerate}
\item \textbf{Connection to the Multivariate Normal Density}: 
The
most important term in the formula \eqref{tdens}  is $(x -
\mu)^T\Sigma^{-1}(x - \mu)$. This exact term also appears in the
multivariate normal density. If $X \sim N(\mu, \Sigma)$, then the
density of $X$ is given by:
\begin{equation*}
  \frac{1}{(2\pi)^{p/2}\sqrt{\det \Sigma}} \exp \left(-\frac{1}{2} (x
    - \mu)^T \Sigma^{-1} (x - \mu) 
  \right). 
\end{equation*}
This suggests that the $t$-density is closely related to the
multivariate normal density. Here is the connection. Suppose $X \sim
N_p(\mu, \Sigma)$ and $V \sim \chi^2_{\nu}$ (this is the chi-squared
distribution with $\nu$ degrees of freedom) are independent. Then 
\begin{equation}\label{t_fact}
  T := \mu + \frac{X - \mu}{\sqrt{V/\nu}} \sim t_p(\mu, \Sigma, \nu). 
\end{equation}
Thus, in the notation $t_{p}(\mu, \Sigma, \nu)$, $\nu$ denotes
degrees of freedom, $p$ denotes dimension, $\mu$ and $\Sigma$ denote
the mean vector and covariance matrix of the corresponding normal
random vector $X$. For completeness, we include a proof of
\eqref{t_fact} in Section \ref{proof_t_norm}.

\item \textbf{Individual Components as well as Linear Combinations of
    Components of $T$ are also $t$-distributed}: Suppose $T \sim
  t_p(\mu, \Sigma, \nu)$ and the components of $T$ are $T_1, \dots,
  T_p$. Then each individual component $T_j$ is also
  $t$-distributed. Also every linear combination $a_0 + a_1 T_1 + a_2
  T_2 +  \dots + a_p T_p$ is also $t$-distributed. To see this, first
  write
  \begin{equation*}
    a_0 + a_1 T_1 + \dots + a_p T_p = a_0 + a^T T 
  \end{equation*}
  where $a$ is the $p \times 1$ vector with components $a_1, \dots,
  a_p$. Using the formula \eqref{t_fact}, we can write
  \begin{equation*}
    a_0 + a^T T = (a_0 + a^T \mu) + \frac{(a_0 + a^TX) - (a_0 +
      a^T\mu)}{\sqrt{V/\nu}} 
  \end{equation*}
  Because $a_0 + a^T X \sim N(a_0 + a^T \mu, a^T \Sigma a)$, the same
  fact \eqref{t_fact} applied to this case gives:
  \begin{equation*}
    a_0 + a^T T \sim t_1(a_0 + a^T \mu, a^T \Sigma a, \nu). 
  \end{equation*}
  In particular, this implies that for each $j = 1, \dots, p$,
  \begin{equation*}
    T_j \sim t_1(\mu_j, \Sigma(j, j), \nu)
  \end{equation*}
  where $\mu_j$ is the $j$th component of $\mu$ and $\Sigma(j, j)$ is
  the $(j, j)$th entry of $\Sigma$. 
 \item \textbf{When $\nu$ is large, $t$ is very close to normal}:
   This can intuitively be seen by noting that when $\nu$ is
large, the term $(x - \mu)^T \Sigma^{-1} (x - \mu)/\nu$ is small so
that
\begin{equation*}
1 + \frac{1}{\nu}  (x - \mu)^T \Sigma^{-1} (x - \mu) \approx \exp
\left(\frac{1}{\nu}  (x - \mu)^T \Sigma^{-1} (x - \mu) \right), 
\end{equation*}
where we used the observation that $1 + z \approx e^z$ when $z$
is small. Thus the $t$-density \eqref{tdens} for large $\nu$ becomes
approximately:
\begin{equation*}
  \exp \left(-\frac{\nu + p}{2 \nu} (x - \mu)^T \Sigma^{-1} (x - \mu)
  \right) \approx \exp \left(-\frac{1}{2} (x - \mu)^T
    \Sigma^{-1} (x - \mu) \right) 
\end{equation*}
because $\frac{\nu + p}{\nu} \approx 1$ when $\nu$ is large. This gets
us the normal density:
\begin{equation*}
  t_{p}(\mu, \Sigma, \nu) \approx N_p(\mu, \Sigma) ~~\text{if $\nu$ is
    large}. 
\end{equation*}
\end{enumerate}

It turns out that \eqref{bayespost_lm} is a special case of
\eqref{tdens} for some $p, \mu, \Sigma, \nu$. To see this, we need to
first rewrite \eqref{bayespost_lm} using matrix notation which we do
in the next section.


\section{Matrix Notation for Multiple Linear Regression}
\begin{equation*}
  y =
  \begin{pmatrix}
    y_1 \\ \cdot \\ \cdot \\ \cdot \\ y_n
  \end{pmatrix} ~~~ X =
  \begin{pmatrix}
    1 & x_{11} & \dots & x_{1m} \\
    1 & x_{21} & \dots & x_{2m} \\
    \cdot & \cdot & \dots & \cdot \\
    \cdot & \cdot & \dots & \cdot \\
    \cdot & \cdot & \dots & \cdot \\
    1 & x_{n1} & \dots & x_{nm}
  \end{pmatrix} ~~~
  \beta =
  \begin{pmatrix}
    \beta_0 \\ \beta_1 \\ \cdot \\ \cdot \\ \cdot \\ \beta_m
  \end{pmatrix} ~~~
  \hat{\beta} =
  \begin{pmatrix}
    \hat{\beta}_0 \\ \hat{\beta}_1 \\ \cdot \\ \cdot \\ \cdot \\ \hat{\beta}_m
  \end{pmatrix}
\end{equation*}
This notation is used not just to write formulae for linear
regression, but also in code. For example, the OLS function in
\texttt{statsmodels} uses the syntax \texttt{sm.OLS(y, X).fit()} to
fit the linear regression model, where $y$ ($n \times 1$ vector) and
$X$ ($n \times (m+1)$ matrix) are defined above. 

With this notation, one can write the sum of squares $S(\beta_0,
\dots, \beta_m)$ as: 
\begin{equation*}
  S(\beta) = S(\beta_0, \dots, \beta_m) = \|y - X \beta\|^2. 
\end{equation*}
There are two important facts about $S(\beta)$:
\begin{enumerate}
\item \textbf{Fact 1}: the least squares estimator $\hat{\beta}$
  is given by the formula:
  \begin{equation}\label{lse}
    \hat{\beta} = (X^T X)^{-1} X^T y. 
  \end{equation}
  The proof of \eqref{lse} is as follows. The gradient of $S(\beta)$
  is given by
  \begin{align*}
    \nabla S(\beta) &= \nabla \left[ \|y - X \beta\|^2 \right] \\
                    &= \nabla \left[ (y - X\beta)^T (y - X \beta) \right] \\
                    &= \nabla \left[ y^T y - \beta^T X^T y - y^T X
                      \beta + \beta^T X^T X \beta\right] = 2 X^T y -
                                                           2 X^T X
                                                           \beta.  
  \end{align*}
  Because $\hat{\beta}$ minimizes $S(\beta)$, the gradient should
  equal zero when $\beta = \hat{\beta}$, and this leads to
  \begin{align}\label{normeq}
    X^T (y - X \hat{\beta}) = 0 \implies X^T X \hat{\beta} = X^T y
    \implies \hat{\beta} = (X^T X)^{-1} X^T y. 
  \end{align}
\item \textbf{Fact 2}: The following Pythagorean identity holds:
  \begin{align}
    \label{pyth}
    S(\beta) = S(\hat{\beta}) + \|X \beta - X \hat{\beta}\|^2 =
    S(\hat{\beta}) + (\beta - \hat{\beta})^T X^T X (\beta - 
    \hat{\beta}). 
  \end{align}
  To prove \eqref{pyth}, write
  \begin{align*}
    S(\beta) &= \|y - X \beta\|^2 \\
             &= \|y - X\hat{\beta} + X \hat{\beta} - X \beta\|^2 \\
    &= \|y - X \hat{\beta}\|^2 + \|X \hat{\beta} - X \beta\|^2 + 2
      \left<y - X \hat{\beta}, X \hat{\beta} - X \beta \right>. 
  \end{align*}
  The cross product is zero (leading to \eqref{pyth}) because:
  \begin{align*}
     \left<y - X \hat{\beta}, X \hat{\beta} - X \beta \right> &=
                                                                \left(X
                                                                \hat{\beta}
                                                                - X \beta
                                                                \right)^T
                                                                \left(y
                                                                - X
                                                                \hat{\beta}
                                                                \right)
    \\
    &= \left( \hat{\beta} - \beta \right)^T X^T (y - X \hat{\beta}) =
      \left( \hat{\beta} - \beta \right)^T \left(X^T y - X^T X
      \hat{\beta} \right) = 0
  \end{align*}
  where we used \eqref{normeq}. 
\end{enumerate}
Using \eqref{pyth}, we can write the posterior density
\eqref{bayespost_lm} as
\begin{align}
  f_{\beta \mid \text{data}}(\beta) &\propto \left(
                                      \frac{S(\hat{\beta})}{S(\beta)}
                                      \right)^{n/2} \nonumber \\
  &= \left(\frac{S(\hat{\beta})}{S(\hat{\beta}) + \left(\beta -
    \hat{\beta} \right)^T X^T X \left(\beta - \hat{\beta} \right)}
    \right)^{n/2} \nonumber \\ &= \left(\frac{1}{1 + \left(\beta -
    \hat{\beta} \right)^T \frac{X^T X}{S(\hat{\beta})} \left(\beta - \hat{\beta} \right)}
    \right)^{n/2}.  \label{post2}
\end{align}
The above formula is a special case of \eqref{tdens} with
\begin{equation*}
x = \beta, ~~~~  p = m+1, ~~~~ \mu = \hat{\beta}, ~~~~ \nu + p = n, ~~~~
\frac{\Sigma^{-1}}{\nu} = \frac{X^T X}{S(\hat{\beta})}
\end{equation*}
or equivalently
\begin{equation*}
  x = \beta, ~~~~  p = m+1, ~~~~ \mu = \hat{\beta}, ~~~~ \nu = n - m -
  1, ~~~~ \Sigma = \frac{S(\hat{\beta})}{n - m - 1} \left(X^T X
  \right)^{-1}. 
\end{equation*}
We thus have
\begin{equation}\label{post3}
  \beta_0, \dots, \beta_m \mid \text{data} \sim
  t_{m+1}\left(\hat{\beta},\frac{S(\hat{\beta})}{n - m - 1} \left(X^T
      X 
  \right)^{-1}, n - m - 1 \right). 
\end{equation}
With the posterior density \eqref{post3}, one can do uncertainty
quantification 
about the parameters $\beta_0, \beta_1, \dots, \beta_m$. One can generate
multiple samples from $t_{m+1}(\hat{\beta}, (S(\hat{\beta})/(n-m-1))
  (X^T X)^{-1}, n-m-1)$ and plot the resulting fitted values to visualize the
  uncertainty in the coefficients.


In \eqref{post3}, the quantity
$S(\hat{\beta})/(n-m-1)$ is the \textbf{frequentist 
  unbiased} estimator for $\sigma^2$, so we denote it by
$\hat{\sigma}^2$:
\begin{equation*}
  \hat{\sigma} := \sqrt{\frac{S(\hat{\beta})}{n-m-1}}. 
\end{equation*}
$\hat{\sigma}$ can also be justified as a Bayesian estimator of
$\sigma$ (this will be a question in Homework three). The terminology
\textbf{Residual Standard Error} is sometimes used for
$\hat{\sigma}$. 

With the notation for $\hat{\sigma}$, the posterior \eqref{post3}
becomes:
\begin{equation}\label{post4}
  \beta_0, \dots, \beta_m \mid \text{data} \sim
  t_{m+1}\left(\hat{\beta},\hat{\sigma}^2 \left(X^T
      X 
  \right)^{-1}, n - m - 1 \right). 
\end{equation}
By one of the facts mentioned about the $t$-distribution, the
posterior of each individual $\beta_j$ is also $t$:  
  \begin{equation}\label{marg4}
    \beta_j \mid \text{data} \sim t_{1}\left(\hat{\beta}_j,
\hat{\sigma}^2 (X^T X)^{j+1, j+1}, n-m-1 \right) 
  \end{equation}
  where $(X^T X)^{j+1, j+1}$ is the $(j+1)^{th}$ diagonal entry 
  of $(X^T X)^{-1}$ (note that we are using the $(j+1)$th diagonal
  entry of $X^TX$ because $\beta_j$ is the $(j+1)$th component of
  $\beta$). Writing this density out, we have
  \begin{equation*}
    f_{\beta_j \mid \text{data}}(\beta_j) \propto \left(\frac{1}{1 +
        \frac{1}{n-m-1} \frac{(\beta_j -
          \hat{\beta}_j)^2}{\hat{\sigma}^2 (X^TX)^{j+1,
            j+1}}}  \right)^{n/2} 
  \end{equation*}
  which implies that
  \begin{equation*}
    \frac{\beta_j - \hat{\beta}_j}{\hat{\sigma} \sqrt{(X^T X)^{j+1,
          j+1}}} \sim \text{univariate standard } t \text{ with }
    n-m-1 \text{ d.f.}
  \end{equation*}
  This can be used to obtain uncertainty intervals for $\beta_j$. If
  $t_{n-m-1, \alpha/2}$ is the point beyond which the $t$-distribution
  (with $n-m-1$ degrees of freedom) assigns probability $\alpha/2$,
  then
  \begin{equation*}
    \P \left\{-t_{n-m-1, \alpha/2} \leq     \frac{\beta_j - \hat{\beta}_j}{\hat{\sigma} \sqrt{(X^T X)^{j+1,
          j+1}}} \leq t_{n-m-1, \alpha/2} \mid \text{data} \right\} =
  1 - \alpha
\end{equation*}
which is same as:
\begin{equation*}
  \P \left\{\hat{\beta}_j - \hat{\sigma} \sqrt{(X^T X)^{j+1,
          j+1}} t_{n-m-1, \alpha/2} \leq \beta_j \leq \hat{\beta}_j - \hat{\sigma} \sqrt{(X^T X)^{j+1,
          j+1}} t_{n-m-1, \alpha/2}  \mid \text{data}  \right\} = 1 -
    \alpha 
\end{equation*}
This interval:
\begin{equation*}
  \left[\hat{\beta}_j - \hat{\sigma} \sqrt{(X^T X)^{j+1,
          j+1}} t_{n-m-1, \alpha/2},  \hat{\beta}_j - \hat{\sigma} \sqrt{(X^T X)^{j+1,
          j+1}} t_{n-m-1, \alpha/2} \right]
  \end{equation*}
  is called the $100(1-\alpha)\%$ Bayesian Credible interval for
  $\beta_j$. It \textbf{exactly coincides} with the frequentist $100(1
  - \alpha)\%$ confidence interval for $\beta_j$.

The degrees of freedom corresponding to the $t$-density in \eqref{post3} is $n - m - 1$ where
$n$ is the number of observations, and $m$ is the number of
covariates. Thus \textbf{if $n-m-1$ is large}, then the posterior
distribution (which is actually $t$) is approximately normal: 
\begin{equation*}
    t_{m+1}\left(\hat{\beta},\frac{S(\hat{\beta})}{n - m - 1} \left(X^T
      X 
  \right)^{-1}, n - m - 1 \right) \approx N_{m+1}(\hat{\beta},
\frac{S(\hat{\beta})}{n - m - 1} \left(X^T X  \right)^{-1}).  
\end{equation*}
In other words, when $n-m-1$ is large, the $t$-density \eqref{post3} is approximately
equal to the $N_{m+1}(\hat{\beta}, \hat{\sigma}^2 (X^T
X)^{-1})$. Further, when $n-m-1$ is large, the distribution \eqref{marg4}
will be close to the normal  distribution $N(\hat{\beta}_j,
\hat{\sigma}^2 (X^T X)^{j+1, j+1})$. The quantity $\hat{\sigma}
\sqrt{(X^T X)^{j+1, j+1}}$ is known as the standard error
corresponding to $\beta_j$.  



\subsection{Proof of \eqref{t_fact}}\label{proof_t_norm}

\begin{proof}[Proof of \eqref{t_fact}]
  Start with the formula: 
\begin{equation*}
  f_T(y) = \int_0^{\infty} f_{T \mid V = x}(y) f_V(x) dx. 
\end{equation*}
Observe that
\begin{equation*}
  T \mid V = x \sim N\left(\mu,  \frac{\nu}{x} \Sigma \right)
\end{equation*}
so that
\begin{align*}
f_{T \mid V = x}(y) &= \frac{1}{(2 \pi)^{p/2}
                      \sqrt{\text{det}(\frac{\nu}{x} \Sigma)}} \exp
                      \left[-\frac{1}{2} (y - \mu)^T \left(\frac{\nu}{x}
                      \Sigma\right)^{-1} (y - \mu) \right] \\
  &= \frac{x^{p/2}}{(2 \pi)^{p/2} \nu^{p/2} \sqrt{\text{det}(\Sigma)}}
    \exp \left(-\frac{x}{2\nu} (y - \mu)^T \Sigma^{-1} (y - \mu)
    \right)
\end{align*}
where we used $\text{det}(\frac{\nu}{x} \Sigma) = (\nu/x)^p
\text{det}(\Sigma)$. As a result
\begin{align*}
  f_T(y) &= \int_0^{\infty} f_{T \mid V = x}(y) f_V(x) dx \\
  &\propto \int_0^{\infty} \frac{x^{p/2}}{(2 \pi)^{p/2} \nu^{p/2}
    \sqrt{\text{det}(\Sigma)}} 
    \exp \left(-\frac{x}{2\nu} (y - \mu)^T \Sigma^{-1} (y - \mu)
    \right) x^{\frac{\nu}{2} - 1} e^{-x/2} dx \\
  &\propto \int_0^{\infty} x^{\frac{p+\nu}{2} - 1} \exp
    \left(-\frac{x}{2} \left[1 + \frac{1}{\nu} (y - \mu)^T \Sigma^{-1}
    (y - \mu)\right] 
    \right) dx. 
\end{align*}
The change of variable
\begin{equation*}
  t = x \left[1 + \frac{1}{\nu} (y - \mu)^T \Sigma^{-1} (y - \mu)\right]
\end{equation*}
leads to
\begin{align*}
  f_T(y) &\propto \frac{1}{\left[1 + \frac{1}{\nu} (y - \mu)^T
  \Sigma^{-1} (y - \mu) \right]^{\frac{\nu + p}{2}}} \int_0^{\infty}
           t^{\frac{\nu+p}{2} - 1} e^{-t/2} dt \\
  &\propto \frac{1}{\left[1 + \frac{1}{\nu} (y - \mu)^T
  \Sigma^{-1} (y - \mu) \right]^{\frac{\nu + p}{2}}}. 
\end{align*}
which proves \eqref{t_fact}. 
\end{proof}





%\begin{thebibliography}{99}

%\bibitem[MacKay and Peto(1995)]{MacKay1995HierarchicalDirichlet}
%D.~J.~C. MacKay and L.~C. Peto,
%\newblock ``A hierarchical Dirichlet language model,''
%\newblock \emph{Natural Language Engineering}, vol.~1,
%no.~3, pp.~289--308, 1995.

%\bibitem[Rubin(1981)]{Rubin1981BayesianBootstrap}
%D.~B. Rubin,
%\newblock ``The Bayesian bootstrap,''
%\newblock \emph{The Annals of Statistics}, vol.~9,
%no.~1, pp.~130--134, 1981.
  
%\bibitem[Fisher(1934)]{Fisher1934}
%R.~A. Fisher,
%\newblock ``Probability, likelihood and quantity of information in the logic of
%uncertain inference,''
%\newblock \emph{Proceedings of the Royal Society of London. Series~A,
%Containing Papers of a Mathematical and Physical Character}, vol.~146,
%no.~856, pp.~1--8, 1934.

%\bibitem[Efron(2010)]{Efron}
%Efron, B. (2010)
%\textit{Large-scale inference: empirical Bayes methods for estimation,
%testing, and prediction}. Cambridge University Press, 2010.
  
%     \bibitem[Jaynes(2003)]{Jaynes} Jaynes, E. T, (2003)
%  \textit{Probability theory: the logic of science}. Cambridge
%  University Press, 2003.

%  \bibitem[Sinai(1992)]{Sinai} Sinai, Y. G, (1992)
%  \textit{Probability theory: an introductory course}. Springer
%  Verlag, 1992.  

%\bibitem[{deGroot(1986)}]{DeGrootBlackwell}DeGroot, M. H. (1986)
%       A conversation with David Blackwell.
%\newblock \emph{Statistical Science}, \textbf{1}, 40--53.

%         \bibitem[Mosteller(1987)]{Mosteller} Mosteller, F, (1987)
%  \textit{Fifty challenging problems in probability with
%    solutions}. Courier Corporation, 1987.

%    \bibitem[MacKay(2003)]{MacKay} MacKay, D. J. C., (2003)
%  \textit{Information theory, inference, and learning
%  algorithms}. Cambridge University Press, 2003.

%\bibitem[{Lindley(2013)}]{Lindley}Lindley, D. V. (2013)
%   \textit{Understanding Uncertainty}. John Wiley and sons, 2013.


%\end{thebibliography}







\end{document}

%%% Local Variables:
%%% mode: latex
%%% TeX-master: t
%%% End:
